\section{Desarrollo de Herramienta en Octave}
\label{chap:codigo_octave}

Para facilitar la verificación de los diagramas de Bode (tanto asíntotas como curvas reales) y el diagrama polar de manera ágil durante las sesiones de laboratorio, se desarrolló un script en el lenguaje Octave. Dicho script permite al usuario ingresar los coeficientes de los polinomios del numerador y denominador de cualquier función de transferencia $H(s)$ y genera automáticamente las gráficas de respuesta en frecuencia.

El código hace uso del paquete \texttt{control} de Octave para definir la función de transferencia y llama a la librería de terceros \texttt{bodas.m} para el trazado preciso de las asíntotas, además de utilizar las funciones nativas para el diagrama de Nyquist (Polar).

El código fuente del script desarrollado se detalla en el Listado \ref{lst:rao_bodes}.

\lstset{
  literate={á}{{\'a}}1 {é}{{\'e}}1 {í}{{\'i}}1 {ó}{{\'o}}1 {ú}{{\'u}}1
           {Á}{{\'A}}1 {É}{{\'E}}1 {Í}{{\'I}}1 {Ó}{{\'O}}1 {Ú}{{\'U}}1
           {ñ}{{\~n}}1 {Ñ}{{\~N}}1 {¡}{{!`}}1 {¿}{{?`}}1
}
\begin{lstlisting}[float=*, language=Matlab, caption={Script de Análisis de Función de Transferencia en Octave (\texttt{rao\_bodes.m}).}, label={lst:rao_bodes}, frame=tb, breaklines=true, showstringspaces=false, basicstyle=\ttfamily\scriptsize, numbers=left, numberstyle=\tiny, keywordstyle=\color{blue}, commentstyle=\color{green!50!black}, stringstyle=\color{red}, xleftmargin=2em, framexleftmargin=1.5em]
% -------------------------------------------------------------------
% Cátedra: Teoría de Circuitos II - UTN FRC
% Script de Análisis de Función de Transferencia Universal
% Utiliza la librería "bodas" para trazar asíntotas exactas.
% -------------------------------------------------------------------

clear; clc; close all;

% Intentamos cargar el paquete de control, necesario para la función tf()
try
    pkg load control;
catch
    warning('El paquete "control" no está instalado o cargado. Puede fallar si no se usa MATLAB.');
end

fprintf('=======================================================\n');
fprintf('  Análisis de H(s) con Trazado Asintótico (Repo: bodas)\n');
fprintf('=======================================================\n\n');

fprintf('Ingrese los polinomios de la Función de Transferencia H(s).\n');
fprintf('Use el formato de vector de coeficientes descendentes en "s".\n');
fprintf('Ejemplo: para s^2 + 200s + 6400, ingrese [1 200 6400]\n\n');

% 1. Ingreso de datos por parte del usuario
num = input('Ingrese el vector del NUMERADOR: ');
den = input('Ingrese el vector del DENOMINADOR: ');

% 2. Creación del sistema
H = tf(num, den);

fprintf('\n-------------------------------------------------------\n');
fprintf('FUNCIÓN DE TRANSFERENCIA H(s) INGRESADA:\n');
display(H);
fprintf('-------------------------------------------------------\n\n');

% 3. Generación del Diagrama de Bode con bodas.m
fprintf('Ejecutando bodas.m para generar el Diagrama Asintótico...\n');
try
    % La función bodas(G) tomará el objeto tf y generará las gráficas
    % comparando el Bode real con la aproximación asintótica.
    [G_out, w] = bodas(H);
    fprintf('¡Diagramas de Bode generados con éxito!\n');
catch err
    fprintf('\n[ERROR] No se pudo ejecutar la función bodas().\n');
    fprintf('-> Por favor, verifique que los archivos "bodas.m" y "tight_subplot.m"\n');
    fprintf('   estén descargados en la misma carpeta donde está ejecutando este script.\n');
    fprintf('Mensaje del sistema: %s\n', err.message);
end

% 4. Generación del Diagrama Polar (Nyquist)
fprintf('\nGenerando Diagrama Polar...\n');
% Abrimos una nueva figura para que bodas.m no la sobreescriba
figure('Name', 'Diagrama Polar de H(s)'); 
nyquist(H);
grid on;
title('Diagrama Polar (Locus de H(jw))');

fprintf('\nAnálisis finalizado. Revise las ventanas de figuras emergentes.\n');
fprintf('=======================================================\n');
\end{lstlisting}
